% MAIN SNAPSHOT — LSPV/RMV::TRACECONJ1::2026-10-12T16:57:37Z
\documentclass[11pt]{article}
\usepackage{fontspec}
\usepackage{geometry}
\usepackage{amsmath,amssymb,mathtools}
\usepackage{booktabs,longtable,array,tabularx}
\usepackage{enumitem}
\usepackage[hidelinks]{hyperref}
\geometry{margin=1in}
\setlist{nosep}
\newcommand{\OPEN}{\textsf{OPEN}}
\newcommand{\EXACT}{\textsf{PROVED/EXACT}}
\newcommand{\KNOWN}{\textsf{KNOWN THEOREM}}
\newcommand{\COMP}{\textsf{COMPUTATIONAL}}
\newcommand{\RET}{\textsf{FALSIFIED/RETIRED}}
\newcommand{\CAND}{\textsf{CANDIDATE}}
\newcommand{\same}{\textnormal{\textsc{same\_realization}}}
\title{Universal-Core Synthesis: The Conjunction of Trace Guards Still Does Not Force Semantic Work}
\author{LSPV-5 / MAIN synthesis}
\date{13 October 2026}

\begin{document}
\sloppy
\maketitle

\section*{Executive status}
\textbf{P versus NP remains open.}  This run combines the strongest surviving
trace-side repairs from the newest peer reports and falsifies their conjunction
as a universal semantic bridge.  For every total uniform deterministic decider
$A$ with a known time-constructible bound, there is an extensionally equivalent
wrapper whose same execution has all of the following: a source-injective
data-bearing trace, a complete input-length-only control/address schedule,
intervention-live paths from every input bit to the output, positive charged
work, and the exact one-bit answer.  Nevertheless, the terminal semantic view
can collide on inputs with different compatibility objects $H$.  The full
trace recovers the source only because the wrapper transports it by prefix
parity; recovering $H$ from that trace is a new computation, not evidence that
$A$ paid for $H$.

\paragraph{Strict status.}
\textsf{OBLIVIOUS-LIVE-SOURCE-INJECTIVE-WRAPPER1} is \EXACT{} as an
elementary wrapper theorem, pending independent review.  The finite identity
and hostile-control battery is \COMP{}.  The proposed universal bridge
\textsf{TRACE-GUARD-CONJUNCTION-AS-SEMANTIC-BRIDGE1} is \RET{}.
\textsf{SEMANTIC-WITNESS-OR-DIRECT-CORE-LB1} remains \OPEN{}.

\section*{Recovery preflight}
The recovery prompt named
\texttt{research/main/2026-10-12\_0850\_MAIN\_SYNTHESIS-0023.tex} as the
newest verified LSPV-5 artifact.  A fresh canonical read found the later
verified snapshot
\texttt{research/main/2026-10-12\_1547\_MAIN\_SYNTHESIS-0025.tex}, with both
its required header and final sentinel.  No exact body for an uncommitted
later synthesis was available, so no scientific text was reconstructed and no
duplicate recovery artifact was written.  The stale observation is diagnosed
as repository-state lag, not a current missing-body condition.

\section*{Corpus reviewed}
This synthesis read the current repository protocol, \texttt{main4.tex},
\texttt{main5.tex}, root \texttt{main.tex}, the current task specifications,
the previous two MAIN snapshots, and the newest four specialist reports:
\begin{itemize}
\item LSPV-1 \textsf{CUTGAP1}: a fixed predeclared polynomial checkpoint
schedule can be hidden; adaptive, program-specific, and all-time transfer
remain open;
\item LSPV-2 \textsf{SERIALTRACE1}: a source-injective full trace can be
created by a linear-cost serial copy, while pointwise extraction of $H$
requires a separate run or recomputation;
\item LSPV-3 \textsf{CAUSALCANCEL1}: every source bit can be an
intervention-live ancestor of the output while naturally induced changes
cancel and the terminal view retains only the answer and one parity bit;
\item LSPV-4 \textsf{MODELBRIDGE-0028}: complete fixed control/address
schedules plus the decision bit collide on equal-length inputs with different
residual compatibility.
\end{itemize}
Primary sources were checked first for oblivious simulation, relativization,
algebrization, natural proofs, and model-specific time--space lower bounds.

\section*{Exact model, encoding, and target}
Let $x=\langle F,\alpha\rangle$ be a fixed self-delimiting binary encoding of
a 3-CNF and a partial assignment, and let $N=|x|$.  Define
\[
 H(F,\alpha)=\{\beta:\alpha\cup\beta\models F\},\qquad
 h(F,\alpha)={\bf 1}[H(F,\alpha)\ne\varnothing].
\]
The clean core $\mathcal C_N$ contains all such encoded pairs of length $N$;
it does not exclude Horn, affine/Tseitin, 2SAT-encoded,
bounded-treewidth, product, padded, redundant, or alternate-encoding controls.

Fix a total uniform deterministic decider $A$ and a known time-constructible
bound $p(N)$ on its running time.  The desired universal theorem would have to
quantify over every such $A$ and prove that its \emph{same realized execution}
either exposes an $H$-specific object or incurs a charge that lands in a
model-matched unconditional lower bound.  The observer, extractor,
reconstructor, simulator, and any second run are all charged in $N$.

\section*{Candidate under attack}
The newest reports separately defeat bare source-injectivity, bare path
causality, and bare control-schedule retention.  A natural repair is to demand
them simultaneously, together with positive execution cost:
\begin{quote}
Every polynomial SAT decider has a same-run view that is data-bearing and
source-injective, follows a fully retained control/address schedule, has live
causal paths from the source to the answer, and performs positive charged work;
therefore the execution must have paid for compatibility.
\end{quote}
Call this \textsf{TRACE-GUARD-CONJUNCTION-AS-SEMANTIC-BRIDGE1}.  The next
theorem supplies all its trace-side antecedents without its semantic
conclusion.

\section*{Theorem: OBLIVIOUS-LIVE-SOURCE-INJECTIVE-WRAPPER1}
Let $A$ be as above.  There is a uniform deterministic wrapper $W_A$ such that
for every $x\in\{0,1\}^N$:
\begin{enumerate}
\item $W_A(x)=A(x)$;
\item its mutable trace contains prefix parities
\[
 p_0=0,\qquad p_i=p_{i-1}\oplus x_i\quad(1\le i\le N),
\]
and is therefore source-injective because $x_i=p_{i-1}\oplus p_i$;
\item its simulated control/address schedule depends only on $N$ and $p(N)$;
\item each source bit has an intervention-live path to the output; and
\item it uses $O(N+p(N)^2)$ time and $O(p(N)+N)$ space in a direct fixed-sweep
implementation.  The displayed costs include input scanning, simulation,
cancellation, and output.
\end{enumerate}
Yet the terminal tuple can be chosen as
\[
 (a,p,q,z,y)=(A(x),p_N,p_N,A(x)\oplus p_N,A(x)),
\]
which is determined by only $(A(x),p_N)$ and need not determine $H(x)$.

\paragraph{Construction and proof.}
First scan $x$ and record the prefix parities.  Next simulate each of the at
most $p(N)$ steps of $A$ by a complete left-to-right and right-to-left sweep of
a $p(N)$-cell encoded work zone.  The sweep order is fixed by $N$ and $p(N)$;
data values affect only the encoded configuration, not the retained
control/address schedule.  This elementary oblivious simulation costs
$O(p(N)^2)$ time and $O(p(N))$ work cells.  Finally set
\[
 p=p_N,\quad q=p,
 \quad z=A(x)\oplus p,
 \quad y=z\oplus q=A(x).
\]
Changing a source bit while holding the rest fixed flips $p$.  Intervening on
the $p$ input to $z$ flips $y$, and intervening on $q$ flips $y$; hence both
paths are live under the standard intervention test.  Under natural execution
the two changes cancel, so extensional correctness is exact.  Prefix
differences reconstruct all of $x$, proving source-injectivity.  The resource
bounds follow by summing the scan, fixed-sweep simulation, and constant-cost
cancellation stages. \hfill$\square$

\paragraph{Terminal collision.}
Use fixed-width encodings of the eight satisfiable 3-CNFs that force the eight
assignments to three variables, each unit constraint represented by a repeated
literal 3-clause.  Their $H$ objects are distinct singleton assignments and
all decision bits equal one.  The terminal tuple has only two possibilities
because it additionally retains one parity bit.  Therefore two formulas of
equal encoded length have identical terminal tuples but different $H$ by the
pigeonhole principle.  The finite audit exhibits the collision
$H=000$ versus $H=011$.

\paragraph{What the theorem does not say.}
The theorem does not show that every conceivable trace view is useless.  It
shows that this conjunction of popular guards is insufficient.  A
construction-specific relational transcript, interactive certificate, or
semantic witness may still be useful, but a separate universal transfer must
force that object in every decider and charge it to the same execution.

\section*{Bounded reproducible audit}
The script
\texttt{experiments/review/trace\_guard\_conjunction\_synthesis.py}
enumerates every binary source of length $1$ through $12$ against four
decision functions (zero, first bit, majority, and parity).  It records:
\begin{itemize}
\item $32{,}760$ source runs;
\item $32{,}760$ exact prefix-trace reconstructions;
\item $32{,}760$ exact output checks;
\item $65{,}520$ intervention-liveness checks;
\item $32{,}760$ fixed-schedule equality checks;
\item $80$ named hostile-control checks across 2SAT, Horn, affine/Tseitin,
bounded-treewidth, independent-product, padding, redundant-clause, and
alternate-encoding families, under DEAD, COPY, DELAY, INTERP, and RECOMP
wrappers; and
\item an explicit equal-length terminal collision among eight distinct forced
assignments.
\end{itemize}
The result is stored at
\texttt{results/raw/review/2026-10-13\_trace\_guard\_conjunction.json} with
SHA-256
\texttt{8c1e2ea72a9bd8227e5f2fe45042418c15f450240dd3f0298c0689d43396b44b}.
These are finite identity/property tests.  They do not estimate asymptotic
running time and are not evidence for an exponential SAT constant.

\section*{Duplication / contradiction audit}
The result is not a duplicate of the component reports.  \textsf{SERIALTRACE1}
establishes cheap source-injectivity but not fixed-schedule simulation or
causal liveness.  \textsf{CAUSALCANCEL1} establishes live paths with
cancellation but does not add oblivious schedule retention.  \textsf{MODELBRIDGE-0028}
establishes schedule/decision collisions but retains no data-bearing source
trace.  The present theorem composes all three controls in one execution and
charges every stage.  No contradiction was found: the component theorems are
mutually consistent and their conjunction strengthens the negative result.

\section*{What changed relative to root main}
The prior root already recorded the separate serial-trace,
causal-cancellation, and schedule-collision failures.  This run adds the
strictly stronger fact that combining those repairs does not recover a
semantic bridge.  The proof spine is updated so that future candidates may not
claim novelty merely by conjoining source-injectivity, a fixed schedule,
causal liveness, and positive charge.  The missing arrow is narrowed to an
$H$-specific semantic witness or a direct lower bound on the clean-core
decision problem in the same general machine model.

\section*{Adopted / rejected / retired prior components}
\begin{center}
\begin{tabularx}{\textwidth}{p{0.27\textwidth}Xp{0.16\textwidth}}
\toprule
Component & Disposition & Status\\
\midrule
Prefix-parity source-injective trace & Adopted as cheap transport, not semantic
work & \EXACT{}\\
Causal cancellation gadget & Adopted as a live-path hostile control & \EXACT{}\\
Fixed-sweep control/address schedule & Adopted as the model-matched oblivious
simulation layer & \KNOWN{} / exact implementation\\
Trace-guard conjunction as compatibility exposure & Rejected by the composed
wrapper & \RET{}\\
Semantic-state count $\Rightarrow$ runtime & Rejected; no transition lower
bound follows & \RET{}\\
Same-realization $H$-specific transfer & Retained as the missing arrow & \OPEN{}\\
\bottomrule
\end{tabularx}
\end{center}

\section*{Faithful comparison with actual main4 and main5}
The repository entry points were read directly.  \texttt{main4.tex} is the ARC
achievable-region manuscript and \texttt{main5.tex} is the WRAPPER construction
manuscript.  They state construction-specific controllers, exact finite
audits, and conditional elimination or promise-domain results.  They preserve
all-input acquisition/coverage obligations and label the separation open.
Neither manuscript asserts that every uniform deterministic SAT decider is an
ARC/WRAPPER realization, exposes $H$, retains a full trace, follows a fixed
sweep, or supplies a causal certificate.  The present synthesis therefore
does not promote their structural claims into arbitrary-computation lower
bounds.

\section*{Barrier and SAME\_REALIZATION audit}
\paragraph{Same realization.}
All positive wrapper properties occur in one execution, so this test does not
mix incompatible realizations.  However, decoding $x$ from prefix parity and
then computing $H(x)$ is a new charged procedure.  Its existence does not show
that the wrapped execution computed, retained, or paid for $H$.

\paragraph{Relativization and algebrization.}
The construction relativizes: an oracle machine can be simulated with the
same oracle, while parity transport and cancellation are unchanged.  It
therefore supplies no escape from Baker--Gill--Solovay.  It uses no
arithmetization and claims no escape from the Aaronson--Wigderson
algebrization barrier.

\paragraph{Natural proofs and nonuniformity.}
No large constructive circuit property is supplied.  Consequently the result
neither invokes nor evades Razborov--Rudich.  The target remains uniform
deterministic computation; a lower bound for circuits, branching programs,
proof systems, or a fixed streaming interface cannot be imported without an
explicit model-preserving transfer.  Each candidate decider may have its own
polynomial exponent; the fixed-sweep square remains polynomial under the
counterfactual assumption.

\paragraph{Time--space lower bounds.}
Known SAT time--space tradeoffs are retained only in their stated restricted
models.  The wrapper theorem gives an upper-bound transformation and no
unconditional superpolynomial endpoint for arbitrary uniform machines.

\paragraph{One-bit output and alternate representations.}
The decider owes only one bit.  Compact circuits, branching programs, affine
elimination, bounded-treewidth dynamic programming, padding, or a wholly
different state coding may bypass any named compatibility representation.
The theorem does not infer runtime from the number of possible $H$ values.

\section*{P=NP counterfactual audit}
Assume SAT has a uniform deterministic decider $A$ running in $N^{k_A}$ time.
The wrapper runs in $O(N^{2k_A}+N)$ time and polynomial space.  Prefix-trace
reconstruction is linear, and ordinary self-reduction can compute witnesses
with polynomially many calls to $A$.  Every derived object remains polynomial
under the assumption.  No model-matched unconditional theorem forbids this.
Thus the wrapper produces no contradiction, and declaring one of these costs
superpolynomial would insert the desired separation as a premise.

\section*{Proof spine and missing arrows}
The surviving spine is
\[
\begin{array}{c}
\text{arbitrary uniform polynomial SAT decider }A\\
\Downarrow\quad\text{\OPEN: $H$-specific same-execution transfer}\\
\text{charged semantic witness, interaction, or restricted realization}\\
\Downarrow\quad\text{\OPEN: model-matched unconditional lower bound}\\
\text{contradiction.}
\end{array}
\]
The new theorem removes the intermediate shortcut
\[
\text{source-injective data trace} + \text{fixed schedule} +
\text{live paths} + \text{positive charge}
\Longrightarrow \text{paid compatibility}.
\]
The next admissible candidate must identify a semantic invariant that the
wrapper cannot obtain by input transport, fixed scheduling, cancellation, or
external recomputation.  Alternatively it must prove a direct clean-core
lower bound in the unrestricted uniform deterministic model.  Both routes are
\OPEN{}.

\section*{Stale or overstated claims found}
The following universal readings are \RET{}:
\begin{itemize}
\item full trace source-injectivity proves that the solver paid for every
function of the input;
\item causal ancestry or intervention liveness proves natural information
flow to the answer;
\item a complete control/address schedule plus a decision bit recovers
compatibility;
\item conjoining those properties repairs their individual failures;
\item positive polynomial work can be converted into a superpolynomial
charge without a landing theorem; and
\item a structural, representation, proof-system, circuit, or time--space
lower bound automatically applies to arbitrary deterministic SAT algorithms.
\end{itemize}

\section*{PQ technique ledger and cross-domain transfer}
\begin{center}
\begin{tabularx}{\textwidth}{p{0.28\textwidth}p{0.22\textwidth}X}
\toprule
Technique & Classification & Reason\\
\midrule
Prefix-difference reconstruction & Exact transfer & Algebraic identity within
the same trace.\\
Fixed-sweep oblivious simulation & Exact transfer & Uniform simulation with
all overhead charged.\\
Causal cancellation & Exact transfer & Boolean identities prove output
preservation and intervention liveness.\\
Circuit/proof-system/time--space lower bounds & Analogy candidate & Endpoint
models do not match arbitrary uniform computation.\\
Trace-guard conjunction $\Rightarrow H$ work & Failed transfer & The composed
wrapper satisfies the antecedent while terminal $H$ collisions persist.\\
Quantum/query or communication analogies & Failed transfer at present & No
proved solver-to-model reduction preserves realization and cost.\\
\bottomrule
\end{tabularx}
\end{center}

\section*{Current ledger}
\begin{center}
\begin{tabular}{p{0.62\textwidth}p{0.20\textwidth}}
\toprule
Item & Status\\
\midrule
Oblivious live source-injective wrapper theorem & \EXACT{}\\
Finite conjunction audit & \COMP{}\\
Trace-guard conjunction as semantic bridge & \RET{}\\
main4/main5 construction claims & As stated; no universal transfer\\
$H$-specific same-execution semantic witness & \OPEN{}\\
Model-matched unconditional contradiction & \OPEN{}\\
P versus NP & \OPEN{}\\
\bottomrule
\end{tabular}
\end{center}

\section*{Peer handoff}
LSPV-1 should attack candidate semantic witnesses with fixed and adaptive
checkpoint hiding after the present wrapper.  LSPV-2 should search for a
relational or interactive object that cannot be reduced to source transport
plus a second computation, and should continue charging reconstruction.
LSPV-3 should test any proposed invariant against two-path cancellation,
state recoding, and equivalent wrappers.  LSPV-4 should require a direct
solver-to-endpoint simulation preserving the same realization, uniformity,
input encoding, and all overhead.  LSPV-5 should reject future candidates that
merely conjoin already-defeated trace properties without an $H$-specific
necessity theorem.

\section*{Primary references checked}
\begin{itemize}
\item N. Pippenger and M. J. Fischer, ``Relations among complexity measures,''
Journal of the ACM 26(2), 1979, 361--381,
\url{https://doi.org/10.1145/322123.322138}.
\item T. Baker, J. Gill, and R. Solovay, ``Relativizations of the $P=?NP$
question,'' SIAM Journal on Computing 4(4), 1975, 431--442,
\url{https://doi.org/10.1137/0204037}.
\item S. Aaronson and A. Wigderson, ``Algebrization: a new barrier in
complexity theory,'' ACM Transactions on Computation Theory 1(1), 2009,
\url{https://doi.org/10.1145/1490270.1490272}.
\item A. A. Razborov and S. Rudich, ``Natural proofs,'' Journal of Computer
and System Sciences 55(1), 1997, 24--35,
\url{https://doi.org/10.1006/jcss.1997.1494}.
\item L. Fortnow, R. Lipton, D. van Melkebeek, and A. Viglas,
``Time-space lower bounds for satisfiability,'' Journal of the ACM 52(6),
2005, 835--865, \url{https://doi.org/10.1145/1101821.1101822}.
\end{itemize}

% END_OF_VERIFIED_MAIN::TRACECONJ1
\end{document}
